\section{Institutional Background}
\label{sec:institutional}

\subsection{Health as a Constitutional Right and the Rise of Litigation}
\label{sec:health_right}

Brazil's 1988 Constitution created the Unified Health System (SUS), which maintains a formulary of approved medications subject to cost-effectiveness criteria \citep{castro2019, soares2019}. Courts, however, have adopted a strict literal interpretation of Article~196 (``health is a right of all and a duty of the state''), under which any individual can obtain virtually any medication through litigation---from common analgesics to rare-disease treatments costing over US\$400,000 per year. Access is inexpensive (a prescription and a public defender suffice), and approximately 85\% of first-instance claims are granted \citep{cnj2019judicializacao}.

The result: 96,000 health court orders nationally in 2017, a 913\% increase in S\~ao Paulo since 2008 (Figure~\ref{fig:distribution}). Litigated health spending alone amounts to approximately US\$300~million annually---nearly 5\% of S\~ao Paulo's public health budget.

\begin{figure}[ht]
  \centering
  \includegraphics[width=0.8\textwidth]{fig_00_litigation_map.pdf}
  \caption{Health Litigation Cases per 1,000 Inhabitants Across Municipalities in S\~ao Paulo (2009--2019)}
  \label{fig:distribution}
  \par\smallskip\footnotesize\textit{Source:} Authors' elaboration based on CNJ/INSPER (2019) litigation data and IBGE population estimates (SIDRA table 6579, 2009--2019 average).
\end{figure}

\subsection{Public Procurement under Judicial Pressure}
\label{sec:procurement_pressure}

The S\~ao Paulo State Department of Health (SES/SP) manages health policy through 99 decentralized public buyer units (PBUs) distributed across the state (Figure~\ref{fig:pbu}).

\begin{figure}[ht]
  \centering
  \includegraphics[width=0.8\textwidth]{fig_00b_litigation_pbu_map.pdf}
  \caption{Public Buyer Units (SES/SP)}
  \label{fig:pbu}
  \par\smallskip\footnotesize\textit{Source:} Authors' elaboration based on SES/SP administrative records.
\end{figure}

Federal Law 8,666/1993 requires all government purchases to follow a formal tendering process in three phases: (i) an \textbf{internal phase} where PBUs identify needs, specify items, set quantities and reference prices, choose the auction format, and publish a public notice; (ii) an \textbf{external phase} where firms compete through sealed-bid (\textit{convite}) or multi-round descending auction (\textit{preg\~ao}); and (iii) a \textbf{delivery phase}.

For ordinary purchases, the internal phase typically spans 30 to 180~days. Court orders fundamentally disrupt this process. Almost all court decisions (99.94\%) are enforced as preliminary injunctions with delivery deadlines of 1 to 10~days, compressing the internal phase to roughly one-third of the ordinary timeline. This compression affects procurement in several ways. First, PBUs lose the ability to aggregate demand across units and over time, reducing their bargaining leverage from bulk purchasing. Second, reference prices must be set quickly with limited market research, often resulting in higher reservation prices. Third, approximately 75\% of litigated items are not on the SUS formulary, meaning PBUs have less experience procuring them. Fourth, roughly 65\% of court orders involve ``packages'' combining SUS and non-SUS items, further complicating procurement planning.

\subsection{The Sanction Channel: Administrative vs.\ Litigated Purchases}
\label{sec:sanction_channel}

A critical institutional feature for our identification strategy is the distinction between \textit{litigated} and \textit{administrative} urgent purchases. Both types face identical planning constraints---compressed timelines, small quantities, and diverted budget resources. However, they differ sharply in the consequences of procurement failure.

For litigated purchases, failure to comply with a court order exposes public officials to severe sanctions: substantial fines (sometimes disproportionate to the purchase value), administrative and criminal liability, and seizure or blocking of public funds. These penalties are public information, as the court order must be referenced in the tender notice.

Since 2009, the SES/SP has attempted to mitigate these costs through administrative requests, a mechanism allowing the government to negotiate supply directly with individuals before litigation occurs. Administrative requests undergo evaluation by a scientific committee using evidence-based medicine criteria. Crucially, failure to complete an administrative purchase carries no penalties for public officials. Between 2009 and 2019, administrative requests generated approximately 9,700 purchases, or about 6.5\% of the total from court orders.

This institutional distinction provides a natural comparison group for isolating the ``under the gun'' effect, since administrative and litigated purchases share all planning and execution constraints but differ only in the penalty regime. Figure~\ref{fig:purchase_types} illustrates the three purchase types. The Online Appendix presents complementary geographic evidence on the distribution of administrative demand.

\begin{figure}[ht]
  \centering
  \includegraphics[width=0.8\textwidth]{fig_00c_litigation_types.pdf}
  \caption{Types of Purchases}
  \label{fig:purchase_types}
  \par\smallskip\footnotesize\textit{Source:} Authors' elaboration.
\end{figure}

\subsection{The Bidding Process}
\label{sec:bidding}

All purchases---ordinary and urgent---are conducted through the \textit{Bolsa Eletr\^onica de Compras} (BEC), S\~ao Paulo's electronic procurement platform, mandatory for common goods since 2007. Two competitive procedures are used: sealed-bid tendering (\textit{convite}) and multi-round descending auctions (\textit{preg\~ao}) with a reverse auction phase. The key implication is that all purchase types use the same platform and procedures; the variation we exploit arises exclusively from differences in planning conditions and penalty regimes.
